\section{System Integration and Testing}
\label{sec:system_integration_and_testing}

The successful deployment of the proposed blockchain-based reporting service requires a rigorous, multi-layered testing strategy. Given the immutable nature of blockchain ledgers, the testing methodology prioritizes defect prevention through a structured framework that encompasses the entire system lifecycle, from individual unit verification to system-wide integration.

\subsection{Testing Governance and Roles}
To ensure comprehensive coverage and objective validation, the testing process is governed by a structured allocation of responsibilities among the project team. Specific Quality Assurance (QA) roles have been mapped to development responsibilities to ensure that each subsystem—blockchain, frontend, AI, and backend—receives dedicated scrutiny.

\begin{table}[h!]
\centering
\caption{Allocation of Testing Roles and Responsibilities}
\label{tab:testing_roles}
\begin{tabular}{|l|p{4cm}|p{7cm}|}
\hline
\textbf{Role Designation} & \textbf{Primary Responsibility} & \textbf{Specific Testing Duties \& Scope} \\ \hline
\textbf{Test Lead} & Overall QA Strategy & Defines the Master Test Plan; orchestrates end-to-end integration tests; manages the CI pipeline. \\ \hline
\textbf{Blockchain Engineer} & Smart Contract Validation & Writes unit tests for Solidity contracts; performs gas optimization analysis; conducts security audits. \\ \hline
\textbf{AI Specialist} & Content Moderation & Curates training/testing datasets; evaluates AI model performance (Precision/Recall/F1); validates the AI filtering logic. \\ \hline
\textbf{Frontend Tester} & DApp Experience & Conducts UI/UX testing; verifies wallet integration (MetaMask) and transaction signing flows. \\ \hline
\end{tabular}
\end{table}

\subsection{Data Acquisition and Synthetic Data Generation}
The project employs a dual-strategy for data acquisition: utilizing verifiable real-world datasets for AI training and generating synthetic datasets for blockchain stress analysis.

\begin{itemize}
    \item \textbf{Real-World Datasets:} For AI content moderation, the system utilizes the \textit{Urban Civic Issues Image Dataset (QR4Change)} for pothole and garbage detection, supplemented by the \textit{Road Damage Detection} and \textit{Garbage Classification V2} datasets. These provide a robust mix of positive and negative samples to ensure high classification accuracy.
    \item \textbf{Synthetic Data Generation:} A Python-based module using the \textit{Faker} library will be developed to programmatically create realistic user reports. This allows for large-scale testing of the blockchain's throughput and latency without compromising user privacy or relying on manual entry.
\end{itemize}

\subsection{Testing Levels}
The testing process follows a hierarchical approach:
\begin{itemize}
    \item \textbf{Unit Testing:} Individual modules, such as smart contract functions and AI inference scripts, are tested in isolation using frameworks like Hardhat and PyTest.
    \item \textbf{Integration Testing:} Focuses on the communication between the DApp, the PoA blockchain, and off-chain IPFS storage.
    \item \textbf{User Acceptance Testing (UAT):} A controlled group of simulated users will interact with the DApp to evaluate usability and functional flow within a local governance context.
\end{itemize}

\section{Performance Evaluation}
\label{sec:performance_evaluation}

The performance evaluation aims to quantify the efficiency, scalability, and reliability of the proposed framework within a controlled experimental environment.

\subsection{Evaluation Metrics}
The system will be evaluated based on the following Key Performance Indicators (KPIs):

\begin{itemize}
    \item \textbf{Blockchain Throughput and Latency:} Measured in Transactions Per Second (TPS) and average block confirmation time on the PoA network.
    \item \textbf{Gas Consumption Analysis:} Evaluation of the computational cost for smart contract operations (e.g., report submission, voting) to ensure economic feasibility.
    \item \textbf{AI Moderation Accuracy:} Measured using Precision, Recall, and F1-score to determine the effectiveness of the content filtering mechanism.
    \item \textbf{Storage Efficiency:} The time taken to upload and retrieve media files from IPFS compared to traditional storage methods.
\end{itemize}

\subsection{Evaluation Environment and Protocols}
Evaluations will be conducted on a private Ethereum-compatible network configured with a Proof of Authority (PoA) consensus mechanism. Key evaluation protocols include:

\begin{itemize}
    \item \textbf{Load Testing:} Utilizing the synthetic data generation module to flood the network with transactions at increasing rates to identify the system's breaking point.
    \item \textbf{Consensus Resilience Test:} Verifying the network's ability to maintain block production during validator node failures, ensuring high availability.
    \item \textbf{AI Robustness Test:} Evaluating model performance against "Challenge Sets" containing noisy or rotated images to ensure reliability in diverse real-world conditions.
\end{itemize}
